% !TeX root = ../main.tex
%%% ---------------------------------------------------------------------------
%%% 实验步骤
%%%
%%% 写作提示：步骤用**编号列表**，一步一个动作，写成别人能照着做的形式。
%%% 不要写成流水账（「然后我们又跑了一遍」），也不要把结果混进来。
%%% 控制变量的说明要写清楚——助教主要看你有没有把不该变的量控制住。
%%% ---------------------------------------------------------------------------
\section{实验步骤}
\label{sec:procedure}

\begin{enumerate}
  \item \textbf{固定随机性。} 对每个种子 $s \in \setof{0,1,2,3,4}$ 统一设置
        \lstinline|torch.manual_seed(s)|、\path{numpy.random.seed(s)} 与
        \lstinline|cudnn.deterministic = True|，保证同一种子下三种优化器拿到
        完全相同的初始权重与数据顺序。这是本实验能做对比的前提。

  \item \textbf{划分数据。} 用固定索引从训练集切出 \num{5000} 张验证集，
        三种优化器共用同一划分，避免划分差异混入结果。

  \item \textbf{训练。} 对每种优化器 $\times$ 每个种子，按\cref{tab:hparam}
        的配置训练 \qty{100}{\epoch}，共 $3 \times 5 = 15$ 次运行。
        每个 epoch 结束后在验证集上评估一次，记录训练损失、验证准确率与用时。

  \item \textbf{选点。} 取验证准确率最高的那个 epoch 的权重作为该次运行的
        最终模型。\emph{注意不要用测试集选点}，否则测试结果会偏乐观。

  \item \textbf{测试。} 用第 4 步选出的权重在测试集上评估一次，记录准确率。

  \item \textbf{学习率敏感性。} 对每种优化器，在其基准学习率的
        $\setof{0.1, 0.3, 1, 3, 10}$ 倍处各跑 3 个种子，考察结论是否随
        $\lr$ 改变。

  \item \textbf{数据处理。} 按\cref{eq:mean-std} 与\cref{eq:uncertainty}
        计算各指标的均值与标准不确定度，按\cref{sec:analysis} 的方法做
        显著性判断。
\end{enumerate}

\begin{caution}
  第 1 步和第 4 步是这个实验最容易出错的地方。种子没固定，三种优化器的
  初始权重不同，测出来的差异有一部分来自初始化而非优化器本身；用测试集
  选点则会让所有结论失去意义。
\end{caution}
